
\chapter{Requisitos Funcionales}


En esta sección se definen y describen las características de alto nivel (requisitos funcionales) del sistema que son esenciales para cubrir las necesidades de los usuarios. Para facilitar su lectura y análisis, los requisitos se presentan en forma de lista estructurada.
\\ \\
% El sector turístico en El Puerto de Santa María, Cádiz, es una de las principales fuentes de actividad económica, destacando por su rica oferta de playas, gastronomía, cultura e historia. Para garantizar una gestión eficiente y digitalizada de los servicios turísticos, es necesario implementar un sistema que optimice la administración de reservas, clientes y proveedores. Este apartado describe los requisitos funcionales esenciales para el correcto funcionamiento de un sistema de gestión en una agencia de viajes local, asegurando que la experiencia del cliente sea ágil, segura y personalizada.


% \section*{RF-1. Gestión de Paquetes Turísticos}

% Gestionar la disponibilidad, precios y detalles de los paquetes turísticos ofrecidos por la agencia.


% \begin{enumerate}[label=RF-1.\arabic*]
%     \item Controlar la disponibilidad y estado de los paquetes turísticos ofrecidos.
%     \begin{enumerate}[label=RF-1.1.\arabic*]
%         \item Crear y editar paquetes turísticos con información detallada y precios actualizados.
%         \item Visualizar la disponibilidad de paquetes turísticos en tiempo real.
%         \item Gestionar la disponibilidad de plazas y fechas de los paquetes turísticos.
%     \end{enumerate}
%     \item Personalización de Ofertas y Recomendaciones
%     \begin{enumerate}[label=RF-1.2.\arabic*]
%         \item Segmentar a los clientes en base a sus preferencias y comportamientos de compra.
%         \item Crear ofertas personalizadas para clientes con perfiles específicos.
%         \item Recomendar paquetes turísticos basados en el historial de viajes de los clientes.
%     \end{enumerate}
% \end{enumerate}

% \section*{RF-1.1 Controlar la disponibilidad y estado de los paquetes turísticos ofrecidos.}


% RF-1.1.1 Crear y editar paquetes turísticos con información detallada y precios actualizados.

% RF-1.1.2 Visualizar la disponibilidad de paquetes turísticos en tiempo real.

% RF-1.1.3 Gestionar la disponibilidad de plazas y fechas de los paquetes turísticos.

% \section*{RF-1.2 Personalización de Ofertas y Recomendaciones}

% RF-1.2.1 Segmentar a los clientes en base a sus preferencias y comportamientos de compra.

% RF-1.2.2 Crear ofertas personalizadas para clientes con perfiles específicos.

% RF-1.2.3 Recomendar paquetes turísticos basados en el historial de viajes de los clientes.

% \begin{enumerate}[label=RF-1.1.\arabic*]
%     \item Crear y editar paquetes turísticos con información detallada y precios actualizados.
%     \item Visualizar la disponibilidad de paquetes turísticos en tiempo real.
%     \item Gestionar la disponibilidad de plazas y fechas de los paquetes turísticos.
% \end{enumerate}

% \section*{RF-1.2 Personalización de Ofertas y Recomendaciones}

% \begin{itemize}
%     \item RF-1.2.1 Segmentar a los clientes en base a sus preferencias y comportamientos de compra.
%     \item RF-1.2.2 Crear ofertas personalizadas para clientes con perfiles específicos.
%     \item RF-1.2.3 Recomendar paquetes turísticos basados en el historial de viajes de los clientes.
% \end{itemize}

% \section*{RF-2. Gestión de Clientes y Reservas}

% Administrar la información de los clientes, sus reservas y preferencias, garantizando una atención personalizada y eficiente.

% \section*{RF-2.1 Registrar y administrar clientes con programas de fidelización y descuentos.}

% RF-2.1.1 Crear perfiles de clientes con información personal y preferencias de viaje.

% RF-2.1.2 Asignar descuentos y beneficios a clientes frecuentes o con membresía.

% \section*{RF-2.2 Facilitar la consulta y actualización de datos de los clientes.}

% RF-2.2.1 Permitir a los clientes acceder a su historial de reservas y datos personales.

% RF-2.2.2 Actualizar la información de los clientes de forma segura y sencilla.

% \section*{RF-2.3 Permitir la reserva de paquetes turísticos con pago anticipado o a plazos.}

% RF-2.3.1 Ofrecer opciones de pago seguro y en línea para reservas de paquetes turísticos.

% RF-2.3.2 Permitir el pago a plazos para paquetes turísticos de alto valor.

% \section*{RF-2.4 Enviar notificaciones automáticas sobre confirmaciones, cambios o cancelaciones de reservas.}

% RF-2.4.1 Notificar a los clientes sobre la confirmación de sus reservas.

% RF-2.4.2 Informar a los clientes sobre cambios en los paquetes turísticos reservados.

% RF-2.4.3 Notificar a los clientes sobre cancelaciones de reservas y reembolsos.

% \section*{RF-3. Gestión de Pagos y Facturación}

% Administrar los pagos de los clientes, generar facturas y gestionar reembolsos de forma eficiente y segura. \\


% RF-3.1 Integrar múltiples métodos de pago como tarjetas bancarias, Bizum y transferencias.

% RF-3.2 Generar facturas automáticas y enviarlas a los clientes.

% RF-3.3 Administrar reembolsos y políticas de cancelación según las condiciones establecidas.

% \section*{RF-4. Portal Web y Atención al Cliente}

% Ofrecer una plataforma web interactiva y un servicio de atención al cliente eficiente para mejorar la experiencia de los usuarios. \\

% RF-4.1 Permitir la reserva y compra de paquetes turísticos a través de una plataforma online.

% RF-4.2 Implementar un chat en vivo para asistencia inmediata a los clientes a traves de implementación de IA.

% RF-4.3 Ofrecer una sección de preguntas frecuentes y asistencia guiada para los usuarios.

% \begin{enumerate}[label=\textcolor{red}{RF-\arabic*.}]
%     \item \textbf{Gestión de Paquetes Turísticos} \\ 
%     Gestionar la disponibilidad, precios y detalles de los paquetes turísticos ofrecidos por la agencia.
%     \begin{enumerate}[label=RF-1.\arabic*]
%         \item Controlar la disponibilidad y estado de los paquetes turísticos ofrecidos.
%         \begin{enumerate}[label=RF-1.1.\arabic*]
%             \item Crear y editar paquetes turísticos con información detallada y precios actualizados.
%             \item Visualizar la disponibilidad de paquetes turísticos en tiempo real.
%             \item Gestionar la disponibilidad de plazas y fechas de los paquetes turísticos.
%         \end{enumerate}
%         \item Personalización de Ofertas y Recomendaciones.
%         \begin{enumerate}[label=RF-1.2.\arabic*]
%             \item Segmentar a los clientes en base a sus preferencias y comportamientos de compra.
%             \item Crear ofertas personalizadas para clientes con perfiles específicos.
%             \item Recomendar paquetes turísticos basados en el historial de viajes de los clientes.
%         \end{enumerate}
%     \end{enumerate}

%     \item \textbf{Gestión de Clientes y Reservas}
%     \begin{enumerate}[label=RF-2.\arabic*]
%         \item Registrar y administrar clientes con programas de fidelización y descuentos.\\
%         Administrar la información de los clientes, sus reservas y preferencias, garantizando una atención personalizada y eficiente.
%         \begin{enumerate}[label=RF-2.1.\arabic*]
%             \item Crear perfiles de clientes con información personal y preferencias de viaje.
%             \item Asignar descuentos y beneficios a clientes frecuentes o con membresía.
%         \end{enumerate}
%         \item Facilitar la consulta y actualización de datos de los clientes.
%         \begin{enumerate}[label=RF-2.2.\arabic*]
%             \item Permitir a los clientes acceder a su historial de reservas y datos personales.
%             \item Actualizar la información de los clientes de forma segura y sencilla.
%         \end{enumerate}
%         \item Permitir la reserva de paquetes turísticos con pago anticipado o a plazos.
%         \begin{enumerate}[label=RF-2.3.\arabic*]
%             \item Ofrecer opciones de pago seguro y en línea para reservas de paquetes turísticos.
%             \item Permitir el pago a plazos para paquetes turísticos de alto valor.
%         \end{enumerate}
%         \item Enviar notificaciones automáticas sobre confirmaciones, cambios o cancelaciones de reservas.
%         \begin{enumerate}[label=RF-2.4.\arabic*]
%             \item Notificar a los clientes sobre la confirmación de sus reservas.
%             \item Informar a los clientes sobre cambios en los paquetes turísticos reservados.
%             \item Notificar a los clientes sobre cancelaciones de reservas y reembolsos.
%         \end{enumerate}
%     \end{enumerate}

%     \item \textbf{Gestión de Pagos y Facturación}\\
%     Administrar los pagos de los clientes, generar facturas y gestionar reembolsos de forma eficiente y segura. 
%     \begin{enumerate}[label=RF-3.\arabic*]
%         \item Integrar múltiples métodos de pago como tarjetas bancarias, Bizum y transferencias.
%         \item Generar facturas automáticas y enviarlas a los clientes.
%         \item Administrar reembolsos y políticas de cancelación según las condiciones establecidas.
%     \end{enumerate}

%     \item \textbf{Portal Web y Atención al Cliente}\\
%     Ofrecer una plataforma web interactiva y un servicio de atención al cliente eficiente para mejorar la experiencia de los usuarios.
%     \begin{enumerate}[label=RF-4.\arabic*]
%         \item Permitir la reserva y compra de paquetes turísticos a través de una plataforma online.
%         \item Implementar un chat en vivo para asistencia inmediata a los clientes a través de implementación de IA.
%         \item Ofrecer una sección de preguntas frecuentes y asistencia guiada para los usuarios.
%     \end{enumerate}
% \end{enumerate}


\begin{enumerate}[label=\textcolor{red}{RF-\arabic*.}]
    \item \textbf{Gestión de Paquetes Turísticos} \\ 
    Gestionar la disponibilidad, precios y detalles de los paquetes turísticos ofrecidos por la agencia.
    \begin{enumerate}[label=\textcolor{red}{RF-1.\arabic*}]
        \item Controlar la disponibilidad y estado de los paquetes turísticos ofrecidos.
        \begin{enumerate}[label=\textcolor{red}{RF-1.1.\arabic*}]
            \item Crear y editar paquetes turísticos con información detallada y precios actualizados.
            \item Visualizar la disponibilidad de paquetes turísticos en tiempo real.
            \item Gestionar la disponibilidad de plazas y fechas de los paquetes turísticos.
        \end{enumerate}
        \item Personalización de Ofertas y Recomendaciones.
        \begin{enumerate}[label=\textcolor{red}{RF-1.2.\arabic*}]
            \item Segmentar a los clientes en base a sus preferencias y comportamientos de compra.
            \item Crear ofertas personalizadas para clientes con perfiles específicos.
            \item Recomendar paquetes turísticos basados en el historial de viajes de los clientes.
        \end{enumerate}
    \end{enumerate}

    \item \textbf{Gestión de Clientes y Reservas}\\
    Administrar la información de los clientes, sus reservas y preferencias, garantizando una atención personalizada y eficiente.
    \begin{enumerate}[label=\textcolor{red}{RF-2.\arabic*}]
        \item Registrar y administrar clientes con programas de fidelización y descuentos.\\
        \begin{enumerate}[label=\textcolor{red}{RF-2.1.\arabic*}]
            \item Crear perfiles de clientes con información personal y preferencias de viaje.
            \item Asignar descuentos y beneficios a clientes frecuentes o con membresía.
        \end{enumerate}
        \item Facilitar la consulta y actualización de datos de los clientes.
        \begin{enumerate}[label=\textcolor{red}{RF-2.2.\arabic*}]
            \item Permitir a los clientes acceder a su historial de reservas y datos personales.
            \item Actualizar la información de los clientes de forma segura y sencilla.
        \end{enumerate}
        \item Permitir la reserva de paquetes turísticos con pago anticipado o a plazos.
        \begin{enumerate}[label=\textcolor{red}{RF-2.3.\arabic*}]
            \item Ofrecer opciones de pago seguro y en línea para reservas de paquetes turísticos.
            \item Permitir el pago a plazos para paquetes turísticos de alto valor.
        \end{enumerate}
        \item Enviar notificaciones automáticas sobre confirmaciones, cambios o cancelaciones de reservas.
        \begin{enumerate}[label=\textcolor{red}{RF-2.4.\arabic*}]
            \item Notificar a los clientes sobre la confirmación de sus reservas.
            \item Informar a los clientes sobre cambios en los paquetes turísticos reservados.
            \item Notificar a los clientes sobre cancelaciones de reservas y reembolsos.
        \end{enumerate}
    \end{enumerate}

    \item \textbf{Gestión de Pagos y Facturación}\\
    Administrar los pagos de los clientes, generar facturas y gestionar reembolsos de forma eficiente y segura. 
    \begin{enumerate}[label=\textcolor{red}{RF-3.\arabic*}]
        \item Integrar múltiples métodos de pago como tarjetas bancarias, Bizum y transferencias.
        \item Generar facturas automáticas y enviarlas a los clientes.
        \item Administrar reembolsos y políticas de cancelación según las condiciones establecidas.
    \end{enumerate}

    \item \textbf{Portal Web y Atención al Cliente}\\
    Ofrecer una plataforma web interactiva y un servicio de atención al cliente eficiente para mejorar la experiencia de los usuarios.
    \begin{enumerate}[label=\textcolor{red}{RF-4.\arabic*}]
        \item Permitir la reserva y compra de paquetes turísticos a través de una plataforma online.
        \item Implementar un chat en vivo para asistencia inmediata a los clientes a través de implementación de IA.
        \item Ofrecer una sección de preguntas frecuentes y asistencia guiada para los usuarios.
    \end{enumerate}
\end{enumerate}

